\begin{helper}[(Q1) integrated along a segment: strict monotonicity of $A$]
Let $n\ge2$, $p>1$ and $r>0$.  There is a constant $C=C(n,p,r)>0$ such that the
following holds.  Let $\varepsilon\ge0$, let $\gamma_0>0$ and
$\gamma_1\le r\gamma_0$, let $\Omega\subseteq\R^n$, and let
$A\colon\R^n\times\R\times\R^n\to\R^n$ satisfy \noderef{EllipticStructureQ12}
with these data --- \emph{all} of that node's clauses, as that node states them;
they are not re-expanded here, and
$D(\varepsilon)=\{h\in\R^n:\varepsilon+|h|^2\ne0\}$ is the set on which its
differentiability clause and its two bounds are stated.  Then for a.e.\
$x\in\R^n$: if
$x\in\Omega$, then for every $z\in\R$ and all $a,b\in\R^n$,
\[
C\,\gamma_0\bigl(\varepsilon+|a|^2+|b|^2\bigr)^{\frac p2-1}|a-b|^2
\ \le\ \bigl(A(x,z,a)-A(x,z,b),\,a-b\bigr) .
\]
All powers are real powers of nonnegative reals, with the convention
$0^{s}=0$ for $s\ne0$.

\noindent\textbf{On the dependence of $C$, which is $(n,p,r)$ as stated and
$(n,p)$ in fact.}  The statement binds $r$ \emph{before} the existential for
$C$, so what it asserts is a constant depending on $n$, $p$ and $r$; under
pm-0007 that quantifier order is the dependence, and the annotation above says
exactly what the order permits.  The constant the proof below actually produces
depends only on $n$ and $p$ --- it comes from
\noderef{MidpointGradientWeightBound}, which fixes it from $p$ alone, times $n$
--- so the statement is not tight in $r$; binding $r$ after $C$ would be a
strictly stronger statement, and it is not taken because no consumer needs it
and because $r$ is the standing uniformity parameter under which the consumers of
this bound fix their own constants.  The hypothesis $\gamma_1\le r\gamma_0$ is
likewise not needed for the conclusion and is retained for the same reason;
retaining it only weakens the statement.  Condition (Q2) of the thesis itself
\emph{is} used, but only to make the segment integral below converge.
\end{helper}
\begin{proof}
Write $s=\frac p2-1$ and let $C=C(p)>0$ be the constant supplied by
\noderef{MidpointGradientWeightBound} for this $p$ (available since $p>1$),
which that node chooses before any dimension; since $n\ge2$ gives $n\ge1$, its
conclusion applies in the present dimension, so that for every
$\varepsilon\ge0$ and all $a,b\in\R^n$
\[
C\bigl(\varepsilon+|a|^2+|b|^2\bigr)^{s}
\ \le\ \int_0^1(1-t)\bigl(\varepsilon+|ta+(1-t)b|^2\bigr)^{s}dt .
\tag{$\dagger$}
\]
This choice is made before $\varepsilon,\gamma_0,\gamma_1,\Omega,A$ and the
data, which is the required quantifier order, and $C>0$.  A constant depending
on $p$ alone is in particular an admissible value for the $C=C(n,p)$ the
present statement asks for.  The hypothesis $n\ge2$ enters only through the
dimension-dependent factor in Step 1 below and, as just noted, through the
$n\ge1$ that \noderef{MidpointGradientWeightBound} requires.

\emph{Choice of $x$.}  By \noderef{EllipticStructureQ12}, (Q1) of the thesis holds
for a.e.\ $x$ and (Q2) of the thesis holds for a.e.\ $x$; let $N_1$ and $N_2$ be the
corresponding null sets and put $N=N_1\cup N_2$, again a null set.  For every
$x\notin N$ with $x\in\Omega$, both (Q1) of the thesis and (Q2) of the thesis are
available at $x$ for all $z$, $h$, $\xi$, $i$, $j$.  The complement
$\R^n\setminus N$ has full Lebesgue measure and is in particular nonempty,
since Lebesgue measure on $\R^n$ is not the zero measure; so this really does
identify points at which the two conditions may be read off.  Because the
conclusion is itself an ``a.e.\ $x$'' statement, it suffices to prove it for
each such $x$.  Fix $x\notin N$ with $x\in\Omega$, fix $z\in\R$ and fix
$a,b\in\R^n$.  Set
\[
\xi=a-b,\qquad \varphi(t)=ta+(1-t)b=b+t\xi,\qquad
w(t)=\bigl(\varepsilon+|\varphi(t)|^2\bigr)^{s} .
\]

\emph{Step 1: the operator bound coming from (Q2) of the thesis.}  Let
$g\in D(\varepsilon)$ and let $DA(x,z,g)$ be the differential of
$A(x,z,\cdot)$ at $g$, which exists because $D(\varepsilon)$ is open and
$A(x,z,\cdot)$ is $C^1$ there by \noderef{EllipticStructureQ12}, since
$x\in\Omega$; (Q2) of the thesis is available at $g$ for the same reason.  Expanding $\xi=\sum_j\xi^je_j$ and using linearity,
$\bigl(DA(x,z,g)\xi\bigr)^i=\sum_{j}\xi^j\,D_{h_j}A^i(x,z,g)$, so by
(Q2) of the thesis and Cauchy--Schwarz (Preamble, applied to
$(|\xi^1|,\dots,|\xi^n|)$ against $(1,\dots,1)$),
\[
\bigl|\bigl(DA(x,z,g)\xi\bigr)^i\bigr|
\le\gamma_1\bigl(\varepsilon+|g|^2\bigr)^{s}\sum_{j=1}^n|\xi^j|
\le\sqrt n\,\gamma_1\bigl(\varepsilon+|g|^2\bigr)^{s}|\xi| ,
\]
and summing the squares over $i$,
\[
\bigl|DA(x,z,g)\xi\bigr|\ \le\ n\,\gamma_1\bigl(\varepsilon+|g|^2\bigr)^{s}|\xi| .
\tag{$*$}
\]

\emph{Step 2: the scalar function along the segment.}  Define
$\psi\colon\R\to\R$ by
\[
\psi(t)=\bigl(A(x,z,\varphi(t)),\,\xi\bigr)
=\sum_{i=1}^nA^i(x,z,\varphi(t))\,\xi^i .
\]
If $a=b$ the assertion is trivial: both sides vanish, the left because
$|a-b|=0$.  So assume $a\ne b$, i.e.\ $\xi\ne0$; then $\varphi$ is injective, so
the set
\[
T=\{t\in[0,1]:\ \varphi(t)\notin D(\varepsilon)\}
=\{t\in[0,1]:\ \varepsilon+|\varphi(t)|^2=0\}
\]
is empty when $\varepsilon>0$ and has at most one element when
$\varepsilon=0$ --- namely the unique $t_0$ with $\varphi(t_0)=0$, if it lies in
$[0,1]$.  In all cases $T$ is a null set, and $[0,1]\setminus T$ is a union of
at most two intervals.

Now $\psi$ is continuous on $[0,1]$, because $\varphi$ is continuous,
$A(x,z,\cdot)$ is continuous on all of $\R^n$ by the first clause of
\noderef{EllipticStructureQ12}, and $v\mapsto(v,\xi)$ is continuous.  And for
$t\in[0,1]\setminus T$ the point $\varphi(t)$ lies in the open set
$D(\varepsilon)$, where $A(x,z,\cdot)$ is $C^1$; since
$t\mapsto\varphi(t)$ is affine with $\varphi'(t)=\xi$ and $v\mapsto(v,\xi)$ is
linear, the chain rule gives that $\psi$ is differentiable at every such $t$
with
\[
\psi'(t)=\bigl(DA(x,z,\varphi(t))\,\xi,\ \xi\bigr)
=\sum_{i,j=1}^nD_{h_j}A^i\bigl(x,z,\varphi(t)\bigr)\,\xi^j\xi^i
\qquad(t\in[0,1]\setminus T).
\]
Note $\psi(1)=(A(x,z,a),\xi)$ and $\psi(0)=(A(x,z,b),\xi)$, so
$\psi(1)-\psi(0)$ is exactly the right-hand side of the assertion.

\emph{Step 3: $\psi'$ is interval-integrable on $[0,1]$.}  Measurability:
$\psi'$ denotes the derivative function of $\psi$, and the derivative of an
arbitrary function $\R\to\R$ is measurable (Preamble item on measurability of
the derivative) --- the item is about an arbitrary function, so it does not
matter that $\psi$ fails to be differentiable at the at most one point of $T$,
where the derivative function takes its default value.  No continuity of the
Jacobian enters this step, which is why that item is the right tool.
Domination: by Cauchy--Schwarz (Preamble) and $(*)$ with $g=\varphi(t)$, which
is legitimate for $t\notin T$ since then $\varphi(t)\in D(\varepsilon)$,
\[
|\psi'(t)|=\bigl|\bigl(DA(x,z,\varphi(t))\xi,\ \xi\bigr)\bigr|
\le\bigl|DA(x,z,\varphi(t))\xi\bigr|\,|\xi|
\le n\,\gamma_1\,|\xi|^2\,w(t)\qquad(t\in[0,1]\setminus T),
\]
hence for almost every $t\in[0,1]$, since $T$ is a null set.
The function $w$ is interval-integrable on $[0,1]$ by
\noderef{SegmentWeightIntegrable}, applied with this $n$, $p$, $\varepsilon$
and these $a,b$; hence so is the constant multiple
$n\gamma_1|\xi|^2w$, and therefore $\psi'$ is interval-integrable on $[0,1]$ by
the Preamble item on domination.

\emph{Step 4: integrating (Q1) of the thesis.}  First,
\[
\bigl(A(x,z,a)-A(x,z,b),\,a-b\bigr)=\psi(1)-\psi(0)=\int_0^1\psi'(t)\,dt .
\]
If $T=\emptyset$ this is the fundamental theorem of calculus (Preamble) in the
form already used, applicable by Steps 2 and 3.  If $T=\{t_0\}$, apply instead
the variant that asks for differentiability on the open interval and continuity
on the closed one (Preamble) on each of $[0,t_0]$ and $[t_0,1]$: $\psi$ is
continuous on $[0,1]$ and differentiable on $(0,t_0)$ and on $(t_0,1)$ by
Step 2, and $\psi'$ is interval-integrable on each subinterval by Step 3, so
\[
\psi(t_0)-\psi(0)=\int_0^{t_0}\psi'(t)\,dt,
\qquad
\psi(1)-\psi(t_0)=\int_{t_0}^1\psi'(t)\,dt ,
\]
and adding these, with additivity of the interval integral over adjacent
intervals for an interval-integrable function, gives the display.  (If
$t_0\in\{0,1\}$ one of the two subintervals is degenerate and the single
remaining application suffices.)

Next, condition (Q1) of the thesis applied at $x$ with $h=\varphi(t)$ and the vector
$\xi$ reads
\[
\gamma_0\bigl(\varepsilon+|\varphi(t)|^2\bigr)^{s}|\xi|^2
=\gamma_0\,w(t)\,|\xi|^2\ \le\
\sum_{i,j}D_{h_j}A^i\bigl(x,z,\varphi(t)\bigr)\xi^i\xi^j=\psi'(t)
\]
for every $t\in[0,1]\setminus T$, hence for almost every $t\in[0,1]$, which is
what the hypothesis supplies: $\varphi(t)\in D(\varepsilon)$ for those $t$ and
for no others.  Both $t\mapsto\gamma_0w(t)|\xi|^2$ and $\psi'$ are
interval-integrable on $[0,1]$ (Step 3 and
\noderef{SegmentWeightIntegrable}), so monotonicity of the integral under an
almost-everywhere inequality (Preamble) gives
\[
\int_0^1\psi'(t)\,dt\ \ge\ \gamma_0|\xi|^2\int_0^1w(t)\,dt .
\tag{$**$}
\]

\emph{Step 5: removing the factor $1-t$.}  The bound $(\dagger)$ supplied by
\noderef{MidpointGradientWeightBound} carries the extra factor $1-t$ in its
integrand, so it does not immediately bound $\int_0^1w$; we record the
comparison rather than leave it to the reader.  For every $t\in[0,1]$ we have
$0\le1-t\le1$, and $w(t)\ge0$ because a real power of a nonnegative base is
nonnegative; hence
\[
0\le(1-t)\,w(t)\le w(t)\qquad(t\in[0,1]).
\]
The majorant $w$ is interval-integrable on $[0,1]$ by
\noderef{SegmentWeightIntegrable}, and $t\mapsto(1-t)w(t)$ is measurable
(a continuous function times a measurable one), so it too is
interval-integrable, by the Preamble item on domination.  Monotonicity of the
integral (Preamble) then gives
\[
\int_0^1(1-t)\,w(t)\,dt\ \le\ \int_0^1w(t)\,dt ,
\]
and combining this with $(\dagger)$,
\[
C\bigl(\varepsilon+|a|^2+|b|^2\bigr)^{s}\ \le\ \int_0^1w(t)\,dt .
\]

\emph{Conclusion.}  Chaining Step 4 and Step 5, and using $\gamma_0>0$ and
$|\xi|^2\ge0$,
\[
\bigl(A(x,z,a)-A(x,z,b),\,a-b\bigr)\ \ge\ \gamma_0|\xi|^2\int_0^1w(t)\,dt
\ \ge\ C\,\gamma_0\bigl(\varepsilon+|a|^2+|b|^2\bigr)^{s}|a-b|^2 ,
\]
which is the assertion for this $x$.  Since $x$ was an arbitrary point of
$\R^n\setminus N$, the assertion holds for a.e.\ $x$.
\end{proof}
